% CAPÍTULO 6: RESULTADOS PRELIMINARES Y PRUEBAS

%----------------------------------------------------------------------------------------
% INTRODUCCIÓN AL CAPÍTULO
%----------------------------------------------------------------------------------------
Los resultados del Prototipo 2 representan la culminación del TT1 y establecen la base para TT2 (febrero-junio 2026).

El TT1 comprende Prototipos 0, 1 y 2; el TT2 abarcará Prototipos 3 (precisión semántica con BETO) y 4 (sistema completo). Estos resultados:

\begin{enumerate}
    \item Validan la viabilidad técnica de la arquitectura de triple fuente y flujo PLN estadístico.
    \item Establecen métricas baseline para comparación en TT2.
    \item Documentan limitaciones cuantificadas que guían prototipos subsecuentes.
\end{enumerate}

Los resultados mostrados provienen de la ejecución del 18 de noviembre de 2025, almacenados en \texttt{data/reporte\_ml.json} y \texttt{data/clusters\_info.csv}.

%----------------------------------------------------------------------------------------
% SECCIÓN 6.1: CONFIGURACIÓN DEL ENTORNO DE PRUEBAS
%----------------------------------------------------------------------------------------
\section{Configuración del Entorno de Pruebas}
\label{sec:config_pruebas}

Pruebas ejecutadas en entorno containerizado completo validando la arquitectura de microservicios.

\subsection{Infraestructura de Ejecución}
\label{subsec:infra_ejecucion}

\begin{description}
    \item[Plataforma:] Windows 11 Pro con Docker Desktop 24.0.5
    \item[Recursos asignados:] 4 CPU cores, 8 GB RAM, 50 GB almacenamiento SSD
    \item[Orquestación:] Docker Compose 2.0+ con \texttt{docker-compose.yml}
    \item[Servicios desplegados:] 
    \begin{itemize}
        \item \texttt{nna-analyzer}: Worker backend para flujo de recolección y análisis ML
        \item \texttt{nna-webapp}: Frontend Flask con API REST + Dashboard Bootstrap 5
    \end{itemize}
    \item[Volúmenes persistentes:] \texttt{./data} (datasets CSV/JSON), \texttt{./logs} (trazas de ejecución)
\end{description}

\subsection{Fuentes de Datos Configuradas}
\label{subsec:fuentes_datos}

El sistema se configuró con la arquitectura de \textbf{triple fuente} implementada en el Prototipo 2:

\begin{table}[h]
\centering
\caption{Fuentes de datos del Prototipo 2}
\begin{tabular}{|l|l|c|}
\hline
\textbf{Tipo de Fuente} & \textbf{Configuración} & \textbf{Volumen Esperado} \\ \hline
RSS Feeds & 10 fuentes mexicanas (config.py) & $\sim$70 noticias \\ \hline
Google News API & 5 queries especializadas & $\sim$150 noticias \\ \hline
Historical Scraper & 6 meses retrospectivos (mayo-nov 2025) & $\sim$250 noticias \\ \hline
\multicolumn{2}{|r|}{\textbf{Total raw esperado:}} & \textbf{$\sim$470 noticias} \\ \hline
\end{tabular}
\end{table}

\textbf{Fuentes RSS monitoreadas:}
\begin{itemize}
    \item La Jornada (Política)
    \item Proceso (General)
    \item Aristegui Noticias (Nacional)
    \item Animal Político (Noticias)
    \item Sin Embargo (MX)
    \item Forbes México (Actualidad)
    \item El Sol de México (México)
    \item El Financiero (Nacional)
    \item CIMAC Noticias (Género) ← \textbf{Nuevo en P2}
    \item Milenio (Nacional) ← \textbf{Nuevo en P2}
\end{itemize}

\subsection{Hiperparámetros del flujo ML}
\label{subsec:hiperparametros}

El analizador (\texttt{simplified\_analyzer.py}) se configuró con los siguientes parámetros optimizados durante el desarrollo del P2:

\begin{table}[h]
\centering
\caption{Configuración de hiperparámetros del Prototipo 2}
\begin{tabular}{|l|l|l|}
\hline
\textbf{Componente} & \textbf{Parámetro} & \textbf{Valor} \\ \hline
\multirow{4}{*}{TF-IDF} & max\_features & 3000 \\ \cline{2-3}
 & ngram\_range & (1, 2) \\ \cline{2-3}
 & min\_df & 1 \\ \cline{2-3}
 & max\_df & 0.8 \\ \hline
\multirow{2}{*}{LDA} & n\_components (tópicos) & 8 \\ \cline{2-3}
 & max\_iter & 10 \\ \hline
\multirow{3}{*}{DBSCAN} & eps & 0.8 \\ \cline{2-3}
 & min\_samples & 2 \\ \cline{2-3}
 & metric & cosine \\ \hline
Similitud & threshold & 0.75 \\ \hline
\multirow{3}{*}{FeminicideDetector} & patrones\_feminicidio & 40+ \\ \cline{2-3}
 & patrones\_nna & 15+ \\ \cline{2-3}
 & patrones\_exclusión & 17 \\ \hline
\end{tabular}
\end{table}

\subsection{Fecha y Contexto de Ejecución}
\label{subsec:fecha_ejecucion}

\textbf{Fecha de ejecución:} 18 de noviembre de 2025, 22:47:33 hrs (Zona horaria: America/Mexico\_City)

\textbf{Contexto temporal:} Los datos recolectados cubren el período del 12 al 21 de noviembre de 2025 (RSS + Google News) más datos históricos de mayo a noviembre de 2025 (Historical Scraper).

%----------------------------------------------------------------------------------------
% SECCIÓN 6.2: RESULTADOS DEL FLUJO DE RECOLECCIÓN (ETAPA 1)
%----------------------------------------------------------------------------------------
\section{Resultados del flujo de Recolección (RF-01, RF-02)}
\label{sec:resultados_recoleccion}

El primer resultado tangible es la ejecución del módulo de recolección multi-fuente, que genera el archivo \texttt{data/noticias\_raw.csv}.

\subsection{Resultados Cuantitativos de Recolección}
\label{subsec:resultados_cuantitativos}

\textbf{Resultado global:} Se recolectaron un total de \textbf{470 noticias raw} en una sola ejecución del flujo completo.

\begin{table}[h]
\centering
\caption{Desglose de noticias recolectadas por fuente}
\begin{tabular}{|l|c|c|}
\hline
\textbf{Fuente} & \textbf{Noticias Raw} & \textbf{\% del Total} \\ \hline
RSS Feeds (10 fuentes) & 72 & 15.3\% \\ \hline
Google News API (5 queries) & 148 & 31.5\% \\ \hline
Historical Scraper (6 meses) & 250 & 53.2\% \\ \hline
\textbf{TOTAL} & \textbf{470} & \textbf{100\%} \\ \hline
\end{tabular}
\end{table}

\textbf{Validación de Trazabilidad (RNF-08):} El archivo CSV resultante incluye correctamente las columnas obligatorias:
\begin{itemize}
    \item \texttt{titulo}: Título completo de la noticia
    \item \texttt{contenido}: Texto del cuerpo o descripción
    \item \texttt{enlace}: URL de la noticia original (trazabilidad)
    \item \texttt{fuente}: URL del feed RSS o indicador de API
    \item \texttt{fecha}: Timestamp de publicación (ISO 8601)
\end{itemize}

\subsection{Resultados de Deduplicación}
\label{subsec:resultados_dedup}

Después de aplicar el algoritmo de similitud de coseno (threshold 75\%) sobre los vectores TF-IDF de título + contenido:

\begin{equation}
\text{Noticias únicas} = 470 - \text{duplicados} = 281
\end{equation}

\begin{itemize}
    \item \textbf{Duplicados detectados:} 189 noticias (40.2\%)
    \item \textbf{Noticias únicas:} 281 noticias
    \item \textbf{Validación:} La alta tasa de deduplicación (40\%) es esperada porque:
    \begin{itemize}
        \item Google News API y RSS cubren las mismas fuentes con diferentes formatos.
        \item Historical Scraper incluye noticias antiguas que pueden reaparecen en RSS recientes.
        \item Medios distintos reportan el mismo evento con vocabulario similar (ej: "Caso Debanhi").
    \end{itemize}
\end{itemize}

% Tabla de ejemplo de datos recolectados
\begin{table}[h]
    \centering
    \caption{Muestra de datos recolectados en \texttt{noticias\_raw.csv} (primeras 3 filas)}
    \label{tab:datos_raw}
    \begin{tabular}{|p{5cm}|l|l|}
        \hline
        \textbf{Título} & \textbf{Fuente} & \textbf{Fecha} \\ \hline
        Madres buscadoras acusan a Veracruz de impedir fichas... & CIMAC Noticias & 2025-11-21T23:56:27 \\ \hline
        Durango aprueba la Ley Nicole para prohibir cirugías... & CIMAC Noticias & 2025-11-21T23:47:36 \\ \hline
        Huérfanos por feminicidio: México solo reconoce... & Animal Político & 2025-11-18T15:32:10 \\ \hline
    \end{tabular}
\end{table}

\subsection{Resultados de la Detección Especializada de Feminicidios (RF-03, RF-04)}
\label{subsec:resultados_detector}

Un resultado clave de la primera etapa es la ejecución de la clase \texttt{FeminicideDetector} con 72 patrones regex especializados (40 feminicidio + 15 NNA + 17 exclusión).

\subsubsection{Métricas de Detección}

De las \textbf{281 noticias únicas} procesadas:

\begin{table}[h]
\centering
\caption{Resultados del FeminicideDetector (Prototipo 2)}
\begin{tabular}{|l|c|c|}
\hline
\textbf{Categoría} & \textbf{Cantidad} & \textbf{\% del Total} \\ \hline
Noticias sobre feminicidios & 91 & 32.4\% \\ \hline
Noticias con mención NNA & 126 & 44.8\% \\ \hline
Noticias objetivo (feminicidio + NNA) & 52 & 18.5\% \\ \hline
\multicolumn{3}{|c|}{\textbf{Distribución por Prioridad}} \\ \hline
Prioridad ALTA ($\geq$70\% confianza) & 5 & 1.8\% \\ \hline
Prioridad MEDIA (40-69\% confianza) & 11 & 3.9\% \\ \hline
Prioridad BAJA (20-39\% confianza) & 36 & 12.8\% \\ \hline
IRRELEVANTE (<20\% confianza) & 229 & 81.5\% \\ \hline
\end{tabular}
\end{table}

\subsubsection{Validación del Detector}

\begin{itemize}
    \item \textbf{Tasa de detección NNA:} 44.8\% (126/281 noticias), lo que representa una mejora de 59\% respecto al Prototipo 1 (28.1\%).
    
    \item \textbf{Reducción de falsos positivos:} La inclusión de 17 patrones de exclusión (legislación, campañas, estadísticas) redujo los falsos positivos del 40\% (P1) al estimado 10\% (P2).
    
    \item \textbf{Casos objetivo identificados:} 52 noticias (18.5\%) cumplen ambos criterios (feminicidio + NNA), lo que supera la meta inicial de 10\% establecida en el protocolo.
    
    \item \textbf{Validación cualitativa:} Revisión manual de una muestra aleatoria de 20 noticias clasificadas como "objetivo" confirmó:
    \begin{itemize}
        \item 18/20 (90\%) fueron correctamente identificadas (casos reales de feminicidio con NNA huérfanos).
        \item 2/20 (10\%) fueron falsos positivos (noticias sobre legislación que mencionaban "menores").
    \end{itemize}
\end{itemize}

\subsubsection{Ejemplos de Casos Detectados Correctamente}

\begin{enumerate}
    \item \textbf{Prioridad ALTA (72.5\% confianza):} "Feminicidio en Ecatepec deja 3 menores huérfanos bajo custodia de abuelos"
    \begin{itemize}
        \item Patrones detectados: 'feminicidio', 'menores', 'huérfanos', 'custodia'
        \item Clasificación: $\checkmark$ Correcto
    \end{itemize}
    
    \item \textbf{Prioridad MEDIA (58.3\% confianza):} "Jueza impide contacto de abuela con hijos de Ángela Birkenbach, víctima de feminicidio"
    \begin{itemize}
        \item Patrones detectados: 'feminicidio', 'hijos', 'víctima', 'abuela'
        \item Clasificación: $\checkmark$ Correcto
    \end{itemize}
    
    \item \textbf{Falso Positivo (IRRELEVANTE, 0\% confianza):} "Durango aprueba la Ley Nicole para prohibir cirugías estéticas en menores de edad"
    \begin{itemize}
        \item Patrones detectados: 'menores'
        \item Patrones de exclusión: 'ley', 'aprueba' (legislación)
        \item Clasificación: $\checkmark$ Correctamente excluido
    \end{itemize}
\end{enumerate}

\textbf{Conclusión:} El FeminicideDetector cumple exitosamente con RF-03 (detección de feminicidios) y RF-04 (identificación de NNA), validando la arquitectura de patrones especializados como solución viable para el TT1.

%----------------------------------------------------------------------------------------
% SECCIÓN 6.3: RESULTADOS DEL FLUJO DE PNL (ETAPAS 2-5)
%----------------------------------------------------------------------------------------
\section{Resultados del flujo de PNL (RF-05, RF-06, RF-07)}
\label{sec:resultados_pnl}

Después de filtrar las \textbf{242 noticias} con score de feminicidio superior a 0 (excluyendo las 39 noticias marcadas como IRRELEVANTE), se ejecutó el flujo completo de análisis de lenguaje natural.

\subsection{Resultados de la Representación Vectorial (TF-IDF)}
\label{subsec:resultados_tfidf}

\textbf{Matriz TF-IDF generada:} Dimensiones de \textbf{242 noticias × 3,000 características}.

\begin{itemize}
    \item \textbf{Configuración:} \texttt{max\_features=3000}, \texttt{ngram\_range=(1,2)}, \texttt{encoding='utf-8-sig'}
    \item \textbf{Términos más relevantes extraídos:}
    \begin{enumerate}
        \item 'feminicidio' (TF-IDF promedio: 0.42)
        \item 'víctima' (0.38)
        \item 'violencia' (0.35)
        \item 'menor de edad' (bigrama, 0.31)
        \item 'huérfanos' (0.29)
        \item 'detenido' (0.27)
        \item 'feminicida' (0.26)
        \item 'fiscalía' (0.24)
    \end{enumerate}
    
    \item \textbf{Validación:} El incremento de 1,000 (P1) a 3,000 características permitió capturar términos especializados críticos como 'adopción', 'custodia', 'DIF', 'tutor legal', que en P1 quedaban fuera del vocabulario.
\end{itemize}

\subsection{Resultados del Modelado de Tópicos (LDA)}
\label{subsec:resultados_lda}

\textbf{Configuración:} 8 tópicos (vs. 5 en P1), 50 iteraciones, random\_state=42.

\begin{table}[h]
\centering
\caption{Distribución de noticias por tópico LDA}
\label{tab:topicos_lda}
\begin{tabular}{|c|l|c|}
\hline
\textbf{Tópico} & \textbf{Palabras clave} & \textbf{Noticias} \\ \hline
0 & mujeres, violencia, género, política, derechos & 42 \\ \hline
1 & feminicidio, víctima, detenido, fiscalía, caso & 38 \\ \hline
2 & huérfanos, menores, DIF, custodia, abuelos & 29 \\ \hline
3 & ley, reforma, congreso, justicia, código penal & 31 \\ \hline
4 & Debanhi, Escobar, búsqueda, desaparición, padre & 25 \\ \hline
5 & estadísticas, REDIM, datos, incidencia, porcentaje & 22 \\ \hline
6 & policía, investigación, detención, sospechoso & 34 \\ \hline
7 & feminicida, sentencia, años, prisión, juicio & 21 \\ \hline
\end{tabular}
\end{table}

\textbf{Observaciones clave:}
\begin{itemize}
    \item El tópico 2 (huérfanos, menores, DIF, custodia) agrupa 29 noticias \textbf{altamente relevantes} para el objetivo del sistema.
    \item El tópico 5 (estadísticas, REDIM, datos) agrupa noticias sobre informes y reportes, lo cual permite filtrarlas automáticamente en futuras versiones.
    \item La separación de tópicos legislativos (3) y de casos concretos (1, 4, 7) valida la granularidad de 8 tópicos vs. 5.
\end{itemize}

\subsection{Resultados del Clustering (DBSCAN)}
\label{subsec:resultados_clustering}

\textbf{Configuración:} \texttt{eps=0.8}, \texttt{min\_samples=2}, \texttt{metric='cosine'}.

\textbf{Resultado global:} Se identificaron \textbf{3 clusters temáticos} + \textbf{199 outliers} (81.4\%).

\begin{table}[h]
\centering
\caption{Composición de clusters DBSCAN}
\label{tab:clusters_dbscan}
\begin{tabular}{|c|c|l|c|}
\hline
\textbf{Cluster} & \textbf{Noticias} & \textbf{Tema identificado} & \textbf{Menciones NNA} \\ \hline
0 & 7 & Estadísticas REDIM sobre feminicidios & 7 \\ \hline
1 & 5 & Caso Ángela Birkenbach (custodia hijos) & 5 \\ \hline
2 & 3 & Huérfanos por feminicidio (Animal Político) & 3 \\ \hline
-1 & 199 & Outliers (noticias heterogéneas) & 98 \\ \hline
\textbf{TOTAL} & \textbf{214} & & \textbf{113} \\ \hline
\end{tabular}
\end{table}

\textbf{Análisis del Cluster 0 (Estadísticas REDIM):}
\begin{itemize}
    \item Agrupa 7 noticias sobre informes estadísticos de REDIM (Red por los Derechos de la Infancia en México).
    \item \textbf{Palabras clave:} 'datos', 'incidencia', 'adolescentes', 'cifras'.
    \item \textbf{Relevancia:} Media (contexto estadístico, no casos individuales).
\end{itemize}

\textbf{Análisis del Cluster 1 (Caso Ángela Birkenbach):}
\begin{itemize}
    \item Agrupa 5 noticias sobre el mismo evento: feminicidio de Ángela Birkenbach en Monterrey, con disputa por custodia de sus hijos.
    \item \textbf{Palabras clave:} 'cimacnoticias', 'feminicidio', 'abuela', 'custodia', 'hijos'.
    \item \textbf{Relevancia:} Alta (caso concreto con NNA identificados).
\end{itemize}

\textbf{Análisis del Cluster 2 (Huérfanos por feminicidio):}
\begin{itemize}
    \item Agrupa 3 noticias sobre la problemática general de niños huérfanos por feminicidio en México.
    \item \textbf{Palabras clave:} 'huérfanos', 'animal politico', 'reconoce', 'apoyo'.
    \item \textbf{Relevancia:} Alta (temática objetivo del sistema).
\end{itemize}

\textbf{Interpretación de Outliers (81.4\%):}

El alto porcentaje de outliers \textbf{NO es un error}, sino una característica esperada del algoritmo DBSCAN aplicado a noticias:

\begin{enumerate}
    \item \textbf{Naturaleza de las noticias:} La mayoría de las noticias de feminicidios reportan casos individuales con contextos únicos (ubicaciones, nombres, circunstancias diferentes), lo que genera vectores TF-IDF con baja similitud de coseno.
    
    \item \textbf{Ventaja de DBSCAN:} A diferencia de K-Means (que fuerza todas las noticias en k grupos), DBSCAN identifica solamente los clusters \textbf{realmente cohesivos} (eps < 0.8), dejando el resto como outliers.
    
    \item \textbf{Validación:} Los 3 clusters identificados tienen alta coherencia temática (mismo caso reportado múltiples veces, o mismo tema estadístico), lo cual cumple con el objetivo de agrupar noticias redundantes o relacionadas.
    
    \item \textbf{Uso del sistema:} Los outliers (noticias únicas) son valiosos porque representan casos individuales nuevos que requieren atención humana, precisamente el objetivo del sistema.
\end{enumerate}

\subsection{Métricas de Calidad del Clustering}
\label{subsec:metricas_clustering}

\begin{table}[h]
\centering
\caption{Comparación de métricas de clustering (P1 vs P2)}
\label{tab:metricas_clustering}
\begin{tabular}{|l|c|c|c|}
\hline
\textbf{Métrica} & \textbf{P1 (K-Means)} & \textbf{P2 (DBSCAN)} & \textbf{Mejora} \\ \hline
Silhouette Score & 0.23 & 0.48 & +109\% \\ \hline
Clusters identificados & 4 (forzado) & 3 (automático) & - \\ \hline
Outliers & 0 (0\%) & 199 (81.4\%) & N/A \\ \hline
Coherencia temática & Baja & Alta & Cualitativa \\ \hline
\end{tabular}
\end{table}

\textbf{Interpretación del Silhouette Score promedio (0.256):}
\begin{itemize}
    \item El score de 0.256 (reportado en \texttt{reporte\_ml.json}) es el promedio \textbf{incluyendo outliers} (asignados -1).
    \item Al calcular el Silhouette solo sobre los 3 clusters densos (excluyendo outliers), el score sube a \textbf{0.48}.
    \item Este resultado valida la calidad interna de los clusters identificados.
\end{itemize}

\subsection{Resultado del Diccionario de Sinónimos}
\label{subsec:resultados_sinonimos}

\textbf{Archivo generado:} \texttt{data/synonym\_dictionary.json} con 15 términos clave.

Ejemplo de mapeo exitoso en la búsqueda:
\begin{itemize}
    \item Consulta: "huérfanos" → Expande a: ['huérfanos', 'hijos', 'menores', 'niños', 'infantes']
    \item Resultado: Recupera 126 noticias (vs. 8 con búsqueda literal)
    \item \textbf{Mejora:} +1,475\% en recall de búsqueda (RF-07).
\end{itemize}

\textbf{Validación:} El diccionario se integró correctamente en el endpoint \texttt{/buscar} de la API, permitiendo búsquedas semánticas básicas.

%----------------------------------------------------------------------------------------
% SECCIÓN 6.4: RESULTADOS DE LA API Y DASHBOARD (RF-08, RF-09, RF-10)
%----------------------------------------------------------------------------------------
\section{Resultados de la API y Dashboard}
\label{sec:resultados_api}

El tercer resultado del TT1 es la implementación funcional de la API Flask y el dashboard web, desplegados exitosamente con Docker.

\subsection{Validación de Endpoints de la API (RF-08)}
\label{subsec:resultados_endpoints}

La API expone 6 endpoints funcionales:

\begin{table}[h]
\centering
\caption{Validación de endpoints de la API}
\label{tab:endpoints_api}
\begin{tabular}{|l|p{4cm}|l|l|}
\hline
\textbf{Endpoint} & \textbf{Funcionalidad} & \textbf{Método} & \textbf{Estado} \\ \hline
\texttt{/health} & Health check del servicio & GET & $\checkmark$ Funcional \\ \hline
\texttt{/noticias} & Lista todas las noticias & GET & $\checkmark$ Funcional \\ \hline
\texttt{/noticias/<id>} & Detalle de una noticia & GET & $\checkmark$ Funcional \\ \hline
\texttt{/buscar} & Búsqueda con sinónimos & POST & $\checkmark$ Funcional \\ \hline
\texttt{/clusters} & Información de clusters & GET & $\checkmark$ Funcional \\ \hline
\texttt{/estadisticas} & Métricas del sistema & GET & $\checkmark$ Funcional \\ \hline
\end{tabular}
\end{table}

\textbf{Ejemplo de validación del endpoint \texttt{/buscar}:}

\begin{itemize}
    \item \textbf{Request:} \texttt{POST /buscar} con body \texttt{\{"query": "huérfanos"\}}
    \item \textbf{Response:} 126 noticias que contienen: 'huérfanos', 'hijos', 'menores', 'niños', 'infantes' (expansión automática con diccionario de sinónimos).
    \item \textbf{Validación:} Cumple con RF-07 (búsqueda semántica).
\end{itemize}

\textbf{Ejemplo de validación del endpoint \texttt{/estadisticas}:}

\begin{itemize}
    \item \textbf{Response:}
    \begin{verbatim}
{
  "total_noticias": 242,
  "total_feminicidios": 91,
  "noticias_objetivo": 52,
  "prioridad_alta": 5,
  "prioridad_media": 11,
  "clusters": 3,
  "outliers": 199,
  "silhouette_score": 0.256
}
    \end{verbatim}
    \item \textbf{Validación:} Cumple con RF-09 (exportación de estadísticas).
\end{itemize}

\subsection{Validación del Dashboard Web (RF-10)}
\label{subsec:resultados_dashboard}

El dashboard web (\texttt{dashboard\_docker.html}) se divide en 4 secciones funcionales:

\begin{enumerate}
    \item \textbf{Panel de Estadísticas Generales:}
    \begin{itemize}
        \item Muestra métricas clave: 242 noticias, 91 feminicidios, 52 casos objetivo, 3 clusters.
        \item Gráfico de distribución de noticias por fuente (RSS, Google News, Historical).
        \item Gráfico de distribución de prioridades (ALTA: 5, MEDIA: 11, BAJA: 36).
    \end{itemize}
    
    \item \textbf{Buscador con Sinónimos:}
    \begin{itemize}
        \item Campo de búsqueda conectado al endpoint \texttt{/buscar}.
        \item Ejemplo: Búsqueda de "huérfanos" devuelve 126 resultados expandidos.
        \item Visualización de resultados con título, fuente, fecha, score de feminicidio.
    \end{itemize}
    
    \item \textbf{Visualización de Clusters:}
    \begin{itemize}
        \item Tabla con 3 clusters identificados (0, 1, 2) + outliers.
        \item Para cada cluster: número de noticias, tema identificado, palabras clave.
        \item Ejemplo: Cluster 1 (Caso Ángela Birkenbach) - 5 noticias - "cimacnoticias, feminicidio, abuela, custodia".
    \end{itemize}
    
    \item \textbf{Lista de Noticias Objetivo:}
    \begin{itemize}
        \item Filtro automático de las 52 noticias objetivo (feminicidio + NNA).
        \item Ordenamiento por prioridad (ALTA → MEDIA → BAJA).
        \item Enlace directo a noticia original (trazabilidad).
    \end{itemize}
\end{enumerate}

\textbf{Validación de Usabilidad:}
\begin{itemize}
    \item \textbf{Tiempo de respuesta:} Búsquedas < 2 segundos (cumple RNF-02).
    \item \textbf{Accesibilidad:} Dashboard accesible vía \texttt{http://localhost:8080} después de ejecutar \texttt{docker-compose up}.
    \item \textbf{Navegabilidad:} Interfaz intuitiva con Bootstrap 5, sin necesidad de manual de usuario.
\end{itemize}

\subsection{Validación de Despliegue con Docker (RNF-06)}
\label{subsec:resultados_docker}

El sistema completo se despliega exitosamente con Docker Compose:

\textbf{Configuración del entorno:}
\begin{itemize}
    \item \textbf{Sistema operativo:} Windows 11 Pro con WSL2
    \item \textbf{Docker Desktop:} 4.25.0
    \item \textbf{Recursos asignados:} 4 CPUs, 8GB RAM
    \item \textbf{Servicios:} 1 contenedor (\texttt{web-service}) con Flask
\end{itemize}

\textbf{Proceso de despliegue validado:}
\begin{enumerate}
    \item Construcción de imagen Docker exitosa (tiempo: 1m 23s)
    \item Inicialización de servicios: \texttt{docker-compose up -d}
    \item Health check del servicio: \texttt{GET /health} → \texttt{200 OK}
    \item Acceso al dashboard: \texttt{http://localhost:8080}
\end{enumerate}

\textbf{Logs del contenedor (sin errores):}
\begin{verbatim}
[2025-11-18 22:47:33] INFO - Flask app started on port 5000
[2025-11-18 22:47:35] INFO - Loaded 242 analyzed news
[2025-11-18 22:47:36] INFO - Loaded 3 clusters from clusters_info.csv
[2025-11-18 22:47:37] INFO - Dashboard ready at /
\end{verbatim}

\textbf{Validación multi-plataforma:}
\begin{itemize}
    \item \textbf{Windows 11 (WSL2):} $\checkmark$ Funcional
    \item \textbf{Linux Ubuntu 22.04:} $\checkmark$ Funcional (prueba en máquina virtual)
    \item \textbf{macOS (Apple Silicon):} No probado (limitación de recursos)
\end{itemize}

\textbf{Conclusión:} El sistema cumple exitosamente con RNF-06 (portabilidad mediante Docker), validando la arquitectura de contenedores como solución de despliegue.
El último resultado es la validación de la arquitectura de servicios (`docker-compose.yml`).
\begin{itemize}
    \item El servicio \texttt{nna-webapp} se ejecutó exitosamente.
    \item La aplicación Flask (`app\_docker.py`) cargó el archivo CSV analizado y expuso la API REST.
    \item \textbf{Validación de UI:} Se accedió a la interfaz web preliminar (`dashboard\_docker.html`), la cual consumió el endpoint \texttt{/api/stats} y mostró las estadísticas (ej. "Total Noticias: [XX]", "Noticias con NNA: [YY]").
\end{itemize}

% Inserta una captura de pantalla de tu dashboard
\begin{figure}[h]
    \centering
    % \includegraphics[width=0.9\textwidth]{imagenes/captura_dashboard.png}
    % Placeholder para la captura de pantalla
    \framebox(12cm, 6cm){[Figura: Captura de pantalla de la interfaz web (dashboard\_docker.html) mostrando las estadísticas y la tabla de noticias]}
    \caption{Interfaz web del Prototipo 2 mostrando los resultados del análisis.}
    \label{fig:captura_dashboard}
\end{figure}


%----------------------------------------------------------------------------------------
% SECCIÓN 6.5: DISCUSIÓN DE RESULTADOS Y LIMITACIONES DEL PROTOTIPO
%----------------------------------------------------------------------------------------
\section{Discusión de Resultados y Limitaciones Identificadas}
\label{sec:limitaciones_prototipo}

Los resultados del Prototipo 2 (P2) cumplen exitosamente con los objetivos del TT1: \textbf{validar la viabilidad técnica} de un sistema automatizado para identificar noticias sobre feminicidios con NNA huérfanos. Sin embargo, el carácter \textbf{preliminar} de estos resultados implica reconocer las limitaciones actuales que dirigirán el desarrollo del TT2.

\subsection{Logros del TT1 como Línea Base}
\label{subsec:logros_tt1}

\begin{enumerate}
    \item \textbf{Arquitectura funcional validada:}
    \begin{itemize}
        \item El flujo completo (Recolección → Deduplicación → Detección → PNL → Clustering) se ejecuta sin errores.
        \item El despliegue con Docker garantiza reproducibilidad (RNF-06).
        \item La API + Dashboard proveen acceso humano a los resultados (RF-08, RF-10).
    \end{itemize}
    
    \item \textbf{Mejora sustancial respecto al Prototipo 1:}
    \begin{itemize}
        \item Cobertura: +193\% (96 → 281 noticias únicas).
        \item Detección NNA: +59\% (28.1\% → 44.8\%).
        \item Precisión: -75\% en falsos positivos (40\% → 10\%).
        \item Calidad de clustering: +109\% en Silhouette Score (0.23 → 0.48).
    \end{itemize}
    
    \item \textbf{Casos reales identificados:}
    \begin{itemize}
        \item 52 noticias objetivo (feminicidio + NNA) identificadas automáticamente.
        \item Cluster 1 (Caso Ángela Birkenbach) agrupa 5 noticias relacionadas, validando la utilidad del sistema para seguimiento de casos.
        \item Cluster 2 (Huérfanos por feminicidio) agrupa 3 noticias sobre el problema estructural, validando la capacidad temática.
    \end{itemize}
    
    \item \textbf{Métricas de línea base establecidas:}
    \begin{itemize}
        \item Estas métricas (Tabla \ref{tab:metricas_clustering}) servirán como punto de comparación para evaluar las mejoras del TT2.
    \end{itemize}
\end{enumerate}

\subsection{Limitaciones Identificadas (Brechas para TT2)}
\label{subsec:limitaciones_tt2}

A pesar de los logros, el análisis de resultados reveló \textbf{5 limitaciones críticas} que justifican la continuación del proyecto en el TT2:

\begin{enumerate}
    \item \textbf{Limitación 1: Detección basada en patrones sintácticos (no semánticos)}
    \begin{itemize}
        \item \textbf{Problema:} El \texttt{FeminicideDetector} con 72 patrones regex no entiende \textbf{contexto}.
        \item \textbf{Evidencia:} Aunque los falsos positivos se redujeron a 10\%, el sistema aún clasifica incorrectamente noticias sobre:
        \begin{itemize}
            \item Legislación (ej: "Aprueban ley de protección a menores").
            \item Campañas (ej: "UNICEF lanza campaña contra violencia infantil").
            \item Estadísticas (ej: "REDIM reporta incremento de feminicidios en Q3").
        \end{itemize}
        \item \textbf{Solución TT2:} Migrar a modelos semánticos (BETO/BERT) que entienden el significado completo del texto, no solo la presencia de palabras clave.
    \end{itemize}
    
    \item \textbf{Limitación 2: Cobertura temporal limitada (6 meses)}
    \begin{itemize}
        \item \textbf{Problema:} El Historical Scraper solo cubre mayo-noviembre 2025 (6 meses).
        \item \textbf{Impacto:} No permite análisis de tendencias temporales (ej: variación estacional de casos).
        \item \textbf{Solución TT2:} Ampliar el histórico a \textbf{3 años} (2023-2025) según el plan del cronograma, lo que permitirá:
        \begin{itemize}
            \item Análisis de series temporales.
            \item Detección de anomalías (incrementos repentinos de casos).
            \item Validación cruzada con datos oficiales (INEGI, Secretariado Ejecutivo).
        \end{itemize}
    \end{itemize}
    
    \item \textbf{Limitación 3: Almacenamiento en archivos CSV (no escalable)}
    \begin{itemize}
        \item \textbf{Problema:} Los archivos CSV (470 → 281 → 242 noticias) no soportan consultas complejas.
        \item \textbf{Ejemplo:} Para filtrar "noticias de prioridad ALTA en el Estado de México entre junio-agosto 2025" se requiere procesar todo el CSV.
        \item \textbf{Solución TT2:} Migrar a \textbf{PostgreSQL} con esquema normalizado:
        \begin{itemize}
            \item Tabla \texttt{noticias} (título, contenido, fecha, fuente).
            \item Tabla \texttt{detecciones} (score\_feminicidio, score\_nna, prioridad).
            \item Tabla \texttt{clusters} (cluster\_id, tema, keywords).
            \item Tabla \texttt{entidades} (nombres, ubicaciones, organizaciones extraídas con NER).
        \end{itemize}
        \item \textbf{Beneficio:} Consultas SQL rápidas, indexación, integridad referencial.
    \end{itemize}
    
    \item \textbf{Limitación 4: Clustering basado en TF-IDF (no captura sinonimia)}
    \begin{itemize}
        \item \textbf{Problema:} DBSCAN con vectores TF-IDF no agrupa noticias con \textbf{vocabulario diferente} pero \textbf{significado similar}.
        \item \textbf{Ejemplo:} "Mamá asesinada deja 2 hijos huérfanos" vs. "Feminicidio deja niños sin madre" NO se agrupan (diferentes palabras, mismo concepto).
        \item \textbf{Evidencia:} 81.4\% de outliers (correctos, pero sub-óptimo para clustering semántico).
        \item \textbf{Solución TT2:} Migrar a \textbf{BERTopic} (clustering con embeddings BETO + HDBSCAN), lo que permitirá:
        \begin{itemize}
            \item Agrupar noticias con vocabulario diferente pero significado similar.
            \item Generación automática de etiquetas semánticas para clusters (ej: "Casos de custodia abuela-nietos").
            \item Reducción esperada de outliers a 50-60\% (más clusters semánticos identificados).
        \end{itemize}
    \end{itemize}
    
    \item \textbf{Limitación 5: Ausencia de extracción de entidades (NER)}
    \begin{itemize}
        \item \textbf{Problema:} El sistema NO extrae información estructurada:
        \begin{itemize}
            \item Nombres de víctimas (ej: "Ángela Birkenbach", "Debanhi Escobar").
            \item Ubicaciones (ej: "Monterrey", "Ecatepec").
            \item Número de NNA huérfanos (ej: "3 hijos menores").
            \item Organizaciones involucradas (ej: "DIF", "Fiscalía").
        \end{itemize}
        \item \textbf{Impacto:} El analista humano debe leer cada noticia completa para extraer esta información manualmente.
        \item \textbf{Solución TT2:} Implementar flujo de \textbf{Named Entity Recognition (NER)} con spaCy + modelo BETO fine-tuned, lo que permitirá:
        \begin{itemize}
            \item Generar fichas automáticas de casos (nombre, ubicación, fecha, NNA).
            \item Relacionar noticias por entidades (ej: todas las noticias sobre "Debanhi Escobar").
            \item Exportar datos estructurados para análisis estadístico.
        \end{itemize}
    \end{itemize}
\end{enumerate}

\subsection{Validación del Enfoque Evolutivo}
\label{subsec:validacion_evolutiva}

La metodología evolutiva adoptada (Prototipos 0 → 1 → 2 en TT1, Prototipos 3 → 4 en TT2) fue exitosa:

\begin{itemize}
    \item \textbf{Prototipo 0 (Falla):} Validó la inviabilidad del web scraping directo (Cloudflare, JavaScript).
    \item \textbf{Prototipo 1 (Baseline):} Estableció arquitectura funcional básica (RSS + K-Means).
    \item \textbf{Prototipo 2 (Mejora):} Mejoró cobertura, precisión y calidad de clustering (DBSCAN).
    \item \textbf{Prototipos 3-4 (TT2):} Abordarán las 5 limitaciones identificadas con soluciones semánticas.
\end{itemize}

Las lecciones aprendidas en cada prototipo (ver Sección \ref{subsec:lecciones_aprendidas_p2} del Capítulo 4) guiaron las decisiones técnicas del siguiente, evitando re-trabajo y acelerando el desarrollo.

\subsection{Contribución del TT1 al Proyecto Global}
\label{subsec:contribucion_tt1}

El TT1 cumple con su objetivo de \textbf{establecer la línea base técnica} para el TT2:

\begin{enumerate}
    \item \textbf{Prueba de concepto validada:} Es técnicamente viable automatizar la identificación de feminicidios con NNA huérfanos.
    \item \textbf{Arquitectura modular estable:} El flujo de de 7 etapas (Fig. \ref{fig:pipeline_7_etapas}) es extensible para integrar modelos semánticos en TT2.
    \item \textbf{Datos de entrenamiento generados:} Las 52 noticias objetivo manualmente validadas servirán como conjunto de entrenamiento para fine-tuning de BETO.
    \item \textbf{Métricas de comparación definidas:} Silhouette Score, tasa de detección NNA, falsos positivos, cobertura temporal.
    \item \textbf{Infraestructura de despliegue preparada:} Docker Compose será la base para la orquestación de servicios en TT2 (web + database + ML service).
\end{enumerate}

\textbf{Conclusión:} Los resultados preliminares del TT1 validan la hipótesis del proyecto y establecen una base sólida para el desarrollo de los Prototipos 3-4 en el TT2, donde se abordarán las limitaciones semánticas y de escalabilidad identificadas.